\chapter{怎样判断健康与生物科技新闻}

\begin{kaichang}
周末早上，你刚拿起手机，家族群里弹出一条消息，是长辈转发的文章，标题写着：“重磅！科学家证实：这种常见水果的提取物能杀死 98\% 的癌细胞，赶紧转发给家人！”配图是一盘鲜红的草莓和一张穿白大褂的人盯着显微镜的照片。

你会怎么做？立刻转发？立刻反驳？还是先去把草莓吃个够？

先别急着选。这本书从原子一路讲到进化，给了你一条完整的因果链；而这篇附录要把这条链压缩成一件随身工具——一份核查清单。下次再遇到“震惊体”的健康或生物科技新闻，你不用懂那项研究的全部细节，只要按清单问六个问题，就能在几分钟内判断：这条新闻值不值得信，值得信到什么程度。读完这篇附录，我们会回来拆解这条草莓新闻。
\end{kaichang}

这本书的正文反复讲一件事：科学结论的分量，不取决于谁说得大声，而取决于证据是怎样得来的。新闻里的健康说法也不例外。几乎每一条“研究发现”背后，都对应着一种特定的研究方式，而不同的研究方式，能支持的结论强度天差地别。下面六个检查点，按“先看研究是什么，再看数字怎么说，最后看谁在说话”的顺序排列。每个检查点都配一条虚构但非常典型的标题党新闻，我们一起示范怎样拆。

\section{检查点一：这是什么类型的研究？}

判断一条健康新闻，第一个问题永远是：\textbf{这个结论是从什么实验里长出来的？}同样是“科学家发现”，在培养皿里发现和在十万人身上发现，完全是两个世界。

常见的研究类型可以排成一架从低到高的“证据阶梯”：

\begin{itemize}[itemsep=2pt]
\item \textbf{细胞实验}：在培养皿里把某种物质加到细胞上，看细胞的反应。便宜、快速，适合筛选“值不值得继续研究”。但培养皿里没有肝脏帮你解毒，没有免疫系统，没有血液循环——细胞直接泡在高浓度的物质里，和你吞下它之后肠道吸收的那一丁点儿，剂量和环境都完全不同。
\item \textbf{动物实验}：通常用小鼠。小鼠是完整的生物，比培养皿接近人一步，但小鼠的代谢、寿命、疾病方式和人仍有明显差异。在小鼠身上有效的候选药物，绝大多数后来在人身上失败了。
\item \textbf{观察性研究}：在人身上\emph{只观察、不干预}。比如问卷调查十万人“平时喝不喝红酒”，再追踪十年看谁得了心脏病。它能发现真实人群中的关联，样本可以很大，但它分不清“关联”背后的原因——这是检查点四要细讲的关键。
\item \textbf{随机对照试验}（RCT）：把志愿者\emph{随机}分成两组，一组接受干预（吃真药），一组接受对照（吃外观一样的安慰剂），最好连医生和病人都不知道谁在哪一组（双盲）。随机分组像一台“混淆因素搅拌机”，把年龄、收入、生活习惯统统平均到两组，剩下的差别才有底气归因于干预本身。它是判断“某种干预对人有没有用”的黄金标准。
\item \textbf{系统综述与汇总分析}：把一个问题的\emph{所有}合格研究收集起来，按统一标准筛选、合并分析。单项研究可能走运或倒霉，把许多项放在一起看，结论才稳。
\end{itemize}

阶梯的意义不是“低层研究没价值”——细胞实验是几乎所有新药故事的起点——而是：\textbf{低层研究的结论不能被说成人身上的结论}。一条新闻如果只字不提研究是在什么对象上做的，你就该提高警惕。

现在拆那条草莓新闻。“提取物能杀死 98\% 的癌细胞”——这个说法多半是真的，只不过它来自一项细胞实验：研究人员把浓缩的提取物直接滴在培养皿的癌细胞上，癌细胞死了 98\%。这毫不奇怪。整本书都在告诉你物质的作用取决于剂量和环境（第 30 章的浓度、第 44 章的酶促条件）：培养皿里的提取物浓度，可能相当于你一次吃下几十公斤草莓；而且你的消化系统会把它拆解、肝脏会把它代谢，真正到达肿瘤的分子寥寥无几。顺便说一句，用同样的逻辑，浓盐水、酒精、漂白剂都能“杀死培养皿里 98\% 的癌细胞”——没人会把这当成治癌新闻。正确的读法是：“这种提取物值得做下一步动物实验”，仅此而已。

顺着这个例子再多说一句，因为它能帮你读懂最常见的正规医药新闻。一种候选分子从培养皿走到药房，平均要闯十年以上的关卡：细胞筛选、动物实验，然后是人体临床试验——一期只在几十名志愿者身上验证“安不安全”，二期在几百名患者身上看“有没有效的苗头、剂量多少合适”，三期在几千名患者身上做随机对照，正式回答“有没有效、副作用多大”。所以当你看到“某新药一期临床结果喜人”，它的真实含义是“安全性初步过关”，距离“能治这个病”通常还有五到八年，而且大部分候选药会在中途倒下。标题把“一期通过”写成“攻克绝症”，就是把阶梯的第一级说成了顶端。

\section{检查点二：样本多大，有没有对照和重复？}

第二个问题针对研究的“体格”：\textbf{研究对象有多少个？有没有对照组？结果被别人重复出来过吗？}

样本为什么重要？因为随机波动在小样本里会变成“巨大的规律”。抛硬币 4 次，全正面的概率有十六分之一，一点也不稀奇；抛 4000 次还全正面，才值得开新闻发布会。样本越小，运气越能冒充结论。

对照组为什么重要？因为很多病\emph{自己会好}。普通感冒不吃药，一周左右也会痊愈；头痛、失眠、腰背疼常常时轻时重。如果没有一组“不接受干预”的人作比较，你分不清好转是干预的功劳，还是身体自己的修复能力——别忘了，你的免疫系统本来就是亿万年进化打磨出的防御系统（第 62 章）。还有安慰剂效应：光是“相信自己在治疗”，就能让人的主观感受真实地改善，所以对照组必须吃到以假乱真的安慰剂才算数。

来看一条典型标题：“亲测有效！我们小区 10 个孩子吃了这种益生菌软糖，一冬天 8 个没感冒！”拆开看：样本只有 10 个，还是同一个小区的孩子（生活环境相似，不能代表别人）；没有对照组——你不知道同样 10 个\emph{不吃}软糖的孩子会有几个感冒，也许也是 8 个，因为大多数孩子一冬天本来就不感冒；没有重复——明年换一个冬天、换一群孩子，结果可能就反过来了。这条“新闻”连观察性研究都算不上，只是一段轶事。轶事可以是研究的起点，永远不能是结论。

还有一个专门坑“亲测党”的陷阱，叫\emph{均值回归}：任何上下波动的指标——疼痛、血压、考试成绩——在你感觉“最糟”的时候，下一次测量大概率会自动好转一些，不管你干预没干预。人们恰恰习惯在腰最疼的那天去做理疗、在感冒最重的那天喝“神水”，随后好转，功劳就记在了干预头上。球星登上杂志封面后成绩下滑的“封面魔咒”，也是同一回事：能上封面，说明正好处在状态顶峰，接下来向平常水平回落再正常不过。判断“好转是不是干预带来的”，唯一可靠的办法还是对照组。

\section{检查点三：相对风险和绝对风险，分别是多少？}

通过了前两关，新闻开始引用数字了。第三个问题是：\textbf{它报的是相对风险，还是绝对风险？}这是标题党最爱玩的数字魔术。

典型标题：“研究证实：每天吃这种加工肉制品，患某癌风险飙升 50\%！”听着吓人。但“飙升 50\%”是相对风险——新风险比旧风险多出百分之五十。绝对风险呢？假设这种癌症在普通人群中的发病率是万分之四，每天吃这种肉制品的人群中是万分之六。从万分之四到万分之六，确实多了 50\%；但换算成绝对值，只多了万分之二——一万个每天吃的人里，大约多出 2 个病例。

\begin{qiaoliang}
相对风险和绝对风险的换算只要一步除法。设不吃某物的人群发病率是 $r_0$，吃的人群发病率是 $r_1$：
\[
\text{相对风险} = \frac{r_1}{r_0}, \qquad \text{绝对风险差} = r_1 - r_0 .
\]
上面的例子里：$r_1/r_0 = 6/4 = 1.5$，所以“风险是 1.5 倍”或“增加 50\%”；而 $r_1 - r_0 = 0.0006 - 0.0004 = 0.0002$，即万分之二。同一个事实，前一种说法让人恐慌，后一种说法让人冷静。\textbf{负责任的报道应该两个都给。}反过来也成立：如果某种疾病极其常见，哪怕相对风险只增加 10\%，绝对增加的病例数也可能很可观。所以光报绝对差、隐瞒相对比例，同样可能误导。再问一句“基数是多少”，魔术就穿帮了。
\end{qiaoliang}

这不是说那条新闻的结论可以无视——如果天天吃、吃几十年，万分之二的差距累积起来依然值得认真对待，公共卫生部门关心的是几亿人的万分之二。但作为个人决策，你需要的是完整的数字，而不是被挑选过的那一半。

\section{检查点四：是“相关”，还是“因果”？}

第四个问题是所有检查点里最深的一个：\textbf{研究只报告了两件事“一起出现”，新闻却写成了“一件事导致另一件事”吗？}

观察性研究能发现相关，却几乎不能单独证明因果。原因很简单：如果 $A$ 和 $B$ 总是一起出现，至少有三种可能——

\begin{itemize}[itemsep=2pt]
\item $A$ 导致 $B$；
\item $B$ 导致 $A$（因果方向反了）；
\item 某个看不见的 $C$ 同时导致 $A$ 和 $B$（混淆因素）。
\end{itemize}

经典例子：城市里冰淇淋销量和溺水人数高度相关。是冰淇淋导致溺水吗？当然不是——是第三个因素“天热”同时推高了两者。这就是混淆因素。再给你一个校园版：小学生的鞋码和词汇量高度相关——鞋越大，认识的词越多。买大鞋能变聪明吗？当然不行，$C$ 是“年龄”：年龄同时推高了鞋码和词汇量。这类例子听着好笑，但换成“红酒”“燕窝”“某国学班”，包装一换，无数人就乖乖掏钱了。

来看健康新闻版：“长期追踪十万人发现：常喝红酒的人心脏病更少，红酒护心！”这是真实世界里反复上演过的剧本。拆解：观察性研究确实发现过这种相关。但“常喝少量红酒的人”和“不喝酒的人”在别的方面也不一样——前者往往收入更高、饮食更好、看病更方便，这些才是可能的 $C$。后来分析更严格的研究显示，所谓“护心效果”大都来自这类混淆，酒精本身对健康的净影响是负面的。因果语言（“护心”“导致”“防癌”）需要随机对照试验级别的证据才配得上，观察性研究只配用“与……相关”“在……人群中更常见”。

反过来，因果方向颠倒的陷阱也很常见：“研究发现，早睡的孩子成绩更好——快让孩子早睡！”也许吧。但也可能是家庭更规律的孩子两者兼得；还可能是学习负担轻的孩子两者兼得。相关是因果的\emph{线索}，不是因果的\emph{证据}——区分这两者，正是随机对照试验存在的意义。

\begin{zhengju}
“相关冒充因果”最昂贵的一次教训，发生在激素替代疗法上。20 世纪八九十年代，多项大型观察性研究（如追踪十二万名护士的美国护士健康研究）发现：绝经后服用雌激素的女性，心脏病发病率明显更低。机制上也说得通——雌激素影响血脂代谢。于是“激素护心”成了主流观点，数以百万计的女性开始长期服用激素替代疗法。

但总有人不放心：吃激素的女性本来就更常看医生、更注重健康，这个混淆因素排除了吗？争论之下，美国启动了“妇女健康倡议”随机对照试验：一万六千多名绝经后女性被\emph{随机}分成服药组和安慰剂组。2002 年，试验被提前叫停——结果与观察性研究相反：服药组的心脏病、中风和乳腺癌风险不降反升。处方量随即断崖式下跌。

这个故事值得记住的有两点。第一，观察性研究没有错，它忠实地报告了真实的相关；错的是把相关直接读成因果。第二，科学的自我修正靠的不是某篇雄辩的文章，而是更严密的研究设计——当两种设计的结论打架时，随机对照试验的话更算数。这正是检查点一那架证据阶梯在真实历史里起的作用。
\end{zhengju}

\section{检查点五：“统计显著”是“效果巨大”吗？}

第五个问题针对一个最容易被望文生义的词：\textbf{“统计显著”不等于“效果显著”，更不等于“效果巨大”。}

“统计显著”（通常写作 $p<0.05$）的意思其实是：\emph{如果}干预真的毫无效果，纯靠运气碰出这么大（或更大）差异的概率不到 5\%。它回答的是“这差异像不像运气”，完全不回答“这差异有多大”。而一个差异够不够“显著”，强烈依赖样本量：样本越大，越小的差异也能通过显著性门槛。

典型标题：“万人大型研究证实：这种益智补剂显著提升记忆力（$p<0.05$）！”拆解：一万人参加，记忆测试满分 100 分，补剂组平均 71.2 分，安慰剂组平均 70.8 分——差异 0.4 分。因为样本巨大，这 0.4 分足以“统计显著”；但对你而言，花几百块钱换 0.4 分，值得吗？统计学家把“差异有多大”叫作\emph{效应量}，它和“显不显著”是两回事。负责任的新闻会同时报告效应量（或像检查点三那样报告绝对数字），只喊“显著”不给大小的，多半是因为大小拿不出手。

\begin{tiaozhan}
$p<0.05$ 还有一个隐蔽的漏洞：它意味着即使干预毫无效果，每做 20 次检验，平均也会有 1 次纯靠运气跨过门槛。如果研究人员同时检验几十种食物、几十种疾病、好几种分组方式，却不告诉你检验了多少次，那么“撞出”一两个显著结果几乎是必然的——这叫多重比较问题。更糟的做法是先把数据翻来覆去地试，试出显著结果再倒过来编故事。防范它的办法之一，正是检查点六要讲的“预注册”：先把问题和分析方法锁死，再开始收集数据。跳过本段不影响主线，但记住一句口诀：\textbf{一个孤零零的 $p$ 值，什么也不能保证。}
\end{tiaozhan}

\section{检查点六：谁出的钱，谁做的验证？}

最后一个问题离开数字，回到人：\textbf{这项研究是谁资助的？结论经过独立检验了吗？}

利益冲突不等于造假——制药公司资助的药效研究当然也可能做得严谨——但它是一个明确的警示信号，因为大量元研究发现，资助方的利益与研究结论之间存在系统性的相关。所以正规期刊强制要求作者声明利益冲突，而负责任的记者会把它写进报道。

比“谁出钱”更硬的三道闸门是：

\begin{itemize}[itemsep=2pt]
\item \textbf{预注册}：研究开始前，把研究问题、样本量、分析方法登记在公开数据库里。事后改问题、挑数据，一查登记记录就会露馅。
\item \textbf{同行评议}：论文发表前，由同领域的独立专家审读挑错。它不完美——会漏掉错误，偶尔也会放过平庸之作——但它是最基本的门槛。连同行评议都没经过的“预印本”，只能算半成品。
\item \textbf{独立重复}：别的实验室、用别的人群，得出同样的结论。单项研究永远可能走运、倒霉或出错；能被独立重复的结论，才配写进教科书。科学界真正相信的从来不是“某一篇论文”，而是许多研究指向同一个方向后形成的共识。
\end{itemize}

典型标题：“国际权威期刊发表：喝某品牌成长奶粉的宝宝智商更高！”拆解：查一下作者声明，研究由该品牌资助；查一下样本，只有 60 个宝宝；查一下“智商更高”，差异换算下来不到 2 分；再去查查有没有其他实验室重复过——没有。四连问下来，这条“权威新闻”就只剩下广告了。顺便提醒：伪科学内容常披着“期刊发表”的外衣，有些收费就发、几乎不审稿的“掠夺性期刊”也是期刊，所以还要看期刊本身的声誉，而不只是“发表了”三个字。

\begin{wuqu}
“只要能查到论文，这个说法就可信。”——不对。论文是科学讨论的\emph{入场券}，不是真理的盖章。每年发表的论文数以百万计，其中相当比例后来被发现无法重复、需要修正，甚至被撤稿；撤稿的原因从数据错误到图像造假都有。可信度的真正阶梯是：单篇论文 $\rightarrow$ 被独立重复的论文 $\rightarrow$ 汇总多项研究的系统综述 $\rightarrow$ 跨领域证据共同支持的科学共识。一条新闻如果只引用“一项新研究”，尤其是“刚刚发表、首次发现、颠覆认知”的研究，正确的反应不是相信也不是嘲笑，而是把这条结论放进“待验证”文件夹，等它接受重复的考验。
\end{wuqu}

\section{一张可以带走的核查清单}

把六个检查点浓缩成一张表。下次看到健康或生物科技新闻，从上到下过一遍，每一项都标“通过、存疑、不合格”：

\begin{table}[htbp]
\centering
\begin{tabular}{p{0.30\textwidth}p{0.30\textwidth}p{0.30\textwidth}}
\toprule
\textbf{核查问题} & \textbf{危险信号} & \textbf{通过的样子} \\
\midrule
研究是什么类型？ & 用细胞或动物实验暗示人体疗效 & 人体随机对照试验，或系统综述 \\
样本多大？有对照和重复吗？ & 几十个对象、无对照组、轶事式“亲测” & 样本充足、随机分组、双盲、有独立重复 \\
相对风险还是绝对风险？ & 只报“飙升 $XX\%$”或只报万分之几 & 两种数字都给出，并说明基数 \\
相关还是因果？ & 观察性研究配“导致”“防癌”“护心”等词 & 用词克制，或因果结论有 RCT 支持 \\
显著还是巨大？ & 只喊“统计显著”，不报效果多大 & 报告效应量或绝对改善幅度 \\
谁资助？谁验证？ & 厂商资助、无预注册、无人能重复 & 利益冲突透明、预注册、经同行评议且有独立重复 \\
\bottomrule
\end{tabular}
\end{table}

三个“存疑”以上，这条新闻就该进“待验证”文件夹；六项全过，它也只代表“目前最好的证据这样说”，而不是永恒真理——还记得吗，激素替代疗法的共识也曾被更严密的研究翻案。科学给你的不是免死金牌，而是\emph{当下最不容易错的判断}。

\begin{shiyan}
\textbf{亲手撞一次“统计显著”}（安全，在家可做）。

材料：一枚硬币，纸笔，或一个能记录结果的表格。

步骤：假设有一种“开运姿势”号称能让硬币更容易出正面。你按这个姿势抛硬币 10 次，记录正面次数。重复整个“10 次一组”的过程，共做 20 组（也可以约同学每人做几组，合并数据）。

观察点：
\begin{itemize}[itemsep=1pt]
\item 每一组里，正面出现几次？有没有哪一组出现 8 次甚至 9 次正面？
\item 如果那组“8 正 2 反”是你唯一做的一组，你会怎么写新闻标题？
\item 把 20 组合起来，正面总数占 200 次抛掷的百分之几？
\end{itemize}

要点：一次“8 正 2 反”完全可能纯靠运气出现——概率约 4.4\%，做 20 组平均就会撞上将近一组。这就是“小样本 + 只报告得意的那次”如何制造假新闻的最小模型。安全提示：无；唯一的危险是抛得太投入忘记写作业。
\end{shiyan}

\section{回拆开场那条草莓新闻}

现在兑现承诺，用完整清单拆解“水果提取物杀死 98\% 癌细胞”：

\begin{enumerate}[itemsep=2pt]
\item \textbf{研究类型}：细胞实验。能支持的结论只是“值得继续研究”，新闻却暗示“人能治癌”——\textbf{不合格}。
\item \textbf{样本与对照}：培养皿里的几个培养孔，无所谓人群样本；通常也没有和现有药物的正面对照——\textbf{存疑}。
\item \textbf{相对与绝对}：“98\%”是相对数字，基数是培养皿里的癌细胞，不是你身体里的风险——\textbf{偷换概念}。
\item \textbf{相关与因果}：本条倒没有这个问题，但它预设了“杀死培养皿癌细胞”与“治好病人的癌症”之间的因果链，而这条链在药物研发史上断掉过成千上万次——\textbf{存疑}。
\item \textbf{显著与巨大}：新闻没提剂量。在足够高的浓度下，食盐也能 100\% 杀死培养皿癌细胞——\textbf{信息缺失}。
\item \textbf{利益与验证}：标题敦促你“赶紧转发”，文末往往附着一个购买链接或公众号二维码——\textbf{利益相关}。
\end{enumerate}

六个检查点，五个红灯。结论：这项研究本身也许是诚实的初步工作，但这条新闻是不可信的转述。正确动作：不转发，不恐慌，也不用去嘲笑长辈——你可以把这份清单转给他们。

\section{把清单用在更大的画面上}

这套核查方法不只对付标题党。面对正规媒体的正经报道，它帮你读出分寸：一期临床试验说某种新药“耐受良好”，意思是“安全性初步过关”，离“有效”还有两三期试验、好几年的路；面对“基因编辑”“人造生命”这类生物科技新闻（第 78 到 80 章的话题），它帮你分清“在培养细胞里做到了”“在动物身上做到了”和“可以安全用于人”之间隔着多少年的验证。

最后划一条边界：这份清单帮你判断\emph{新闻里的说法有多可靠}，它不能替代医生。你或家人的具体诊疗问题——吃不吃某种药、做不做某项检查——取决于个体的病史、体质和正在用的其他治疗，请务必交给专业人员，并用你在清单里练出的提问方式去问医生：“这项检查改变治疗方案的概率有多大？”“这个药的绝对获益是多少？”好医生欢迎这样的问题。

而这套清单真正的根基，是你在这 87 章里走完整条因果链后养成的直觉：物质的作用取决于结构与剂量，生命是分子按规则运行的系统，证据的强弱取决于获取它的方式。新闻每天都在变，这些规则不变。

\section*{练一练}

\begin{sikaoti}
\item 找一条真实的健康新闻（或回忆一条家族群里见过的），用六问清单逐项打分，写下每一项你是“通过、存疑还是不合格”，并说明依据。
\item 某新闻说：“每天吃两把坚果的人，心脏病风险降低 30\%。”已知该人群中心脏病的年发病率从 1.0\% 降到 0.7\%。分别算出相对风险和绝对风险差，然后各写一句标题：一句只用相对数字，一句只用绝对数字。哪句更像标题党？为什么两种说法都是“对的”？
\item 某“研究”称：“教室里摆放绿植的班级，学生平均分更高。”请给出三种可能的因果解释（$A\rightarrow B$、$B\rightarrow A$、$C$ 同时导致两者），并设计一个能把它们区分开的随机对照试验：分组怎么做？测量什么？要控制哪些混淆因素？
\item 某补剂广告引用一项 $p<0.05$ 的研究：“效果显著！”请写出你要追问的三个问题，并解释为什么“$p<0.05$”本身回答不了其中任何一个。
\item 为什么“被独立重复”比“发表在著名期刊”更能增加一条结论的可信度？请用检查点五和六里讲到的机制（运气、多重比较、利益冲突）解释。
\item 画一张“证据阶梯”图：从下往上标出细胞实验、动物实验、观察性研究、随机对照试验、系统综述，并在每一级旁边写一句“这一级能说什么”和“这一级不能说什么”。在图旁注明：阶梯表示证据强度的人为约定，不是自然界的结构。
\end{sikaoti}
